\definition{Bioproceso}{Secuencia ordenada de pasos para transformar materia prima (generalmente biológica) en producto usando agentes biológicos y sus partes.}

\section{Tipos de bioprocesos}

Podemos distinguir entre 2 tipos de bioprocesos: reactivos y extractivos.

En el caso de los bioprocesos reactivos, estos constan del uso de agentes biológicos (bacterias, hongos) para generar reacciones bioquímicas de interés.

Estos bioprocesos cuentan con 3 etapas: la primera, la etapa de upstream, consta en la preparación de lo involucrado en la reacción, desde el microorganismo usado como del sutrato de este; la segunda etapa, la de midstream o biorreacción, es la reacción planeada, y es aquí donde se produce el producto de interés; y por último, la etapa de downstream, consiste en la separación del producto de interés del resto de sustancias no deseadas (microorganismos restantes, subproductos no deseados, etc.).

Por otro lado, los procesos extractivos consisten en la extracción de productos ya existentes en el agente biológico de interés. Estos productos son lo que se conoce como metabolito.

Este tipo de bioprocesos, a diferencia de los reactivos, constan de solo 2 etapas: una etapa de upstream, donde se prepara la materia orgánica utilizada, y una de downstream, donde después de la extracción se busca purificar el metabolito de interés.

Como se observar entonces al comparar los dos tipos de bioprocesos, la diferencia radica en la ausencia de una biorreacción en los proesos extractivos, y esto es por la naturaleza misma de tal tipo de bioproceso: solo consta de la obtención de compuestos que ya se encuentran presentes en el organismo.